\clearpage\thispagestyle{otFront}\phantomsection\label{references}\pdfbookmark[0]{Sources et édition}{bm-references}\otPageTitre{Sources et édition}{Références officielles · France}\fontsize{11}{15}\selectfont
\par\medskip\textbf{1. Arrêté et calendrier}\par Arrêté du 10 avril 2025 publié au BO n° 16 du 17 avril 2025 : le programme de mathématiques du cycle 3 s’applique en classe de Sixième depuis la rentrée de l’année scolaire 2025-2026.\par\href{https://www.education.gouv.fr/bo/2025/Hebdo16/MENE2504620A}{\textcolor{otGreenDark}{Consulter la source officielle}}\par
\par\medskip\textbf{2. Programme officiel de mathématiques du cycle 3}\par Rubriques « Sixième » : nombres entiers et décimaux, fractions, algèbre ; grandeurs et mesures ; espace et géométrie ; organisation et gestion de données et probabilités ; proportionnalité ; initiation à la pensée informatique.\par\href{https://www.education.gouv.fr/sites/default/files/programme-de-math-matiques-pour-le-cycle-3-439827.pdf}{\textcolor{otGreenDark}{Consulter la source officielle}}\par
\par\medskip\textbf{3. Accompagnement du programme}\par Éduscol : ressources d’accompagnement du programme de mathématiques au cycle 3, dont les exemples de réussite en Sixième, les schémas en barres et la résolution de problèmes.\par\href{https://eduscol.education.gouv.fr/5712/ressources-d-accompagnement-du-programme-de-mathematiques-au-cycle-3}{\textcolor{otGreenDark}{Consulter la source officielle}}\par
\par\bigskip Sources consultées le 23 octobre 2026.
Les cours sont rédigés pour ce livre ; les exercices, défis et bilans
sont originaux.\par\medskip
Périmètre de la Sixième : entiers jusqu’aux milliards ; écritures d’un
décimal (virgule, fraction décimale, nombre mixte, pourcentage) ; arrondis
à l’unité, au dixième et au centième, y compris de nombres non décimaux ;
produit de deux décimaux ; division euclidienne par un entier inférieur à
100 et division décimale par un entier inférieur à 10 ; fraction quotient,
égalités à trous, graduation d’un segment ; égalités, comparaisons,
sommes et différences de fractions, produit par un entier ; pourcentages ;
schémas en barres et motifs évolutifs, sans lettres formelles ;
périmètre du disque, conversions d’aires entre m$^2$, dm$^2$ et cm$^2$,
aires du carré et du rectangle, cm$^3$ ; durées en base soixante et en
heures décimales ; distances, milieu, cercle et disque, médiatrice,
angles, bissectrice, triangles, somme des angles, cercle circonscrit ;
symétrie axiale sur quadrillage, papier pointé et papier uni ; assemblages
de cubes ; enquêtes, tableaux et filtrage ; probabilités entre 0 et 1 et
fréquences ; instructions, séquences, entrées, sorties et répétitions.\par\medskip
La proportionnalité est traitée dans le cadre des grandeurs, par la
linéarité et le retour à l’unité, avec tableaux et échelles ; conformément
au programme, le produit en croix n’est pas enseigné. L’ordre des
chapitres est une proposition de progression.

\clearpage
